\documentclass[UTF8,zihao=-4]{ctexart}
\usepackage[a4paper,margin=2.2cm]{geometry}
\usepackage{fontspec}
\usepackage{tabularx}
\usepackage{array}
\usepackage{ragged2e}
\usepackage{fancyhdr}
\usepackage{hyperref}
\setCJKmainfont{Noto Serif SC}
\setCJKsansfont{Noto Sans SC}
\setmainfont{Noto Serif SC}
\setlength{\parindent}{0pt}
\setlength{\parskip}{0.45em}
\pagestyle{fancy}
\fancyhf{}
\fancyhead[L]{商务智能}
\fancyfoot[C]{\thepage}
\hypersetup{colorlinks=true,linkcolor=black,urlcolor=blue}
\begin{document}
\title{13、OLAP的评价准则}
\author{}
\date{}
\maketitle

Codd 准则：\par
1. 多维概念视图\par
OLAP 必须能从多个维度考察对象。\par
2. 透明性准则\par
OLAP 在体系结构中的位置和数据源对用户透明。\par
3. 存取能力准则\par
能把 OLAP 概念视图映射到异质数据存储上，并提供单一、完整、连续的用户视图。\par
4. 稳定的报表功能\par
当维数和综合层次增加时，报表能力和响应速度不应明显下降。\par
5. 客户 / 服务器体系结构\par
服务器负责统一概念模式、逻辑模式和物理模式；客户端负责应用逻辑和界面。\par
6. 维的同等性原则\par
每一数据维在数据结构和操作能力上都是等同的。\par
7. 动态稀疏矩阵处理准则\par
OLAP 工具应适应特定模型的数据分布，尤其要处理稀疏数据。\par
8. 多用户支持能力准则\par
支持并发访问、数据完整性和安全性机制。\par
9. 非受限的跨维操作\par
对维之间固有层次关系可以自动推导；其他计算需要提供计算完备的语言。\par
10. 直观的数据操作\par
11. 灵活的报表生成\par
12. 不受限维与聚集层次\par
维数不应小于 15，并支持任意聚集层次。\par

FASMI 准则\par

Fast：快速\par
Analysis：分析\par
Shared：共享\par
Multidimensional：多维\par
Information：信息\par

\end{document}
